\documentclass[../../main.tex]{subfiles}%
\graphicspath{{\subfix{./image/}}{\subfix{../../image/}}}

\begin{document}

\section{参考文献与注记}\label{section:ch03-05}

主要依据同样是 Federer\cite{federer1969} 第 1、3 章。\cref{theorem:ch03-3-1}来自 Simon\cite{simon1984} 第 5.1 节。% XREF p=157 "Theorem 3.1 is from Simon [S, Section 5.1]"
Rademacher 定理（\cref{theorem:ch03-3-2}）的证明——我们取自\cite{simon1984} 第 5.2 节——应归于 Morrey（参看\cite{morrey1966} 第 65 页）。% XREF p=157 "The proof of Rademacher's Theorem, which we took from [S, Section 5.2], is due to Morrey (cf. [My, p. 65])"
\cref{section:ch03-01}中的\cref{theorem:ch03-3-3}即\cite{federer1969} 第 3.2.8 节。关于 Lipschitz 连续函数的微分学，另见 Clarke\cite{clarke1983}。% XREF p=157 "in Section 3.1" | p=157 "Theorem 3.3 in Section 3.1 is [F, Section 3.2.8]"

\cref{section:ch03-02}中关于线性映射与 Jacobi 行列式的讨论主要依据 Hardt\cite{hardt1979}。S. Antman 协助我们完成了极分解定理的证明，A. Damlamian 提供了 Binet--Cauchy 公式的计算。另见 Gantmacher\cite{gantmacher1977} 第 9--12、276--278 页。% XREF p=157 "The discussion of linear maps and Jacobians in Section 3.2 is strongly based on Hardt [H]"

\cref{section:ch03-03}中面积公式的证明源自\cite{federer1969} 第 3.2.2--3.2.5 节，并遵循 Hardt 在\cite{hardt1979} 中的论述。我们在\cref{section:ch03-04}中对余面积公式的证明也紧密遵循\cite{hardt1979}，而后者又基于\cite{federer1969} 第 3.2.8--3.2.13 节。\cref{theorem:ch03-3-14}来自\cite{federer1969} 第 3.2.34 节。% XREF p=157 "The proof of the area formula in Section 3.3" | p=157 "Our proof in Section 3.4 of the coarea formula" | p=157 "Theorem 3.14 is from [F, Section 3.2.34]"

\end{document}
